\section{Preliminaries} \label{sec:prelims}

\subparagraph*{Graphs} We only consider finite, undirected, simple graphs. Such a graph~$G$ consists of a set~$V(G)$ of vertices and a set~$E(G) \subseteq \binom{V(G)}{2}$ of edges. For ease of notation, we write~$uv$ for an undirected edge~$\{u,v\} \in E(G)$; note that~$uv = vu$. When it is clear which graph is referenced from context, we write $n$ and $m$ to denote the number of vertices and edges in this graph respectively. For a vertex~$v \in V(G)$, its open neighborhood is~$N_G(v) := \{u \in V(G) \mid uv \in E(G)\}$ and its closed neighborhood is~$N_G[v] := N_G(v) \cup \{v\}$. For a vertex set~$S \subseteq V(G)$ we define its open neighborhood as~$N_G(S) := (\bigcup _{v \in S} N_G(v)) \setminus S$ and its closed neighborhood as~$N_G[S] := \bigcup_{v \in S} N_G[v]$. The subgraph of~$G$ induced by a vertex set~$S \subseteq V(G)$ is the graph~$G[S]$ on vertex set~$S$ with edges~$\{uv \in E(G) \mid \{u,v\} \subseteq S\}$. We use~$G-S$ as a shorthand for~$G[V(G) \setminus S]$ and write~$G - v$ instead of~$G - \{v\}$ for singletons. A \emph{walk} is a sequence of (not necessarily distinct) vertices~$(v_1, \ldots, v_k)$ such that~$v_i, v_{i+1} \in E(G)$ for each~$i \in [k-1]$. The walk is \emph{closed} if we additionally have~$v_k, v_1 \in E(G)$. A \emph{cycle} is a closed walk whose vertices are all distinct. The \emph{length} of a cycle~$(v_1, \ldots, v_k)$ is~$k$. A \emph{path} is a walk whose vertices are all distinct. The \emph{length} of a path~$(v_1, \ldots, v_k)$ is~$k-1$. The vertices~$v_1, v_k$ are the \emph{endpoints} of the path. For two (not necessarily disjoint) vertex sets~$S,T$ of a graph~$G$, we say that a path~$P = (v_1, \ldots, v_k)$ in~$G$ is an~$(S,T)$-path if~$v_1 \in S$ and~$v_k \in T$. If one (or both) of $S$ and $T$ contains only one element, we may write this single element instead of the singleton set consisting of it. 

The \emph{parity} of a path or cycle refers to the parity of its length. For a walk~$W = (v_1, \ldots, v_k)$, we refer to its vertex set as~$V(W) = \{v_1, \ldots, v_k\}$. Observe that if~$W$ is a closed walk of odd parity (a \emph{closed odd walk}), then the graph~$G[V(W)]$ contains a cycle of odd length (an \emph{odd cycle}): any edge connecting two vertices of~$V(W)$ that are not consecutive on~$W$ splits the walk into two closed subwalks, one of which has odd length.

For a positive integer~$q$, a \emph{proper $q$-coloring} of a graph~$G$ is a function~$f \colon V(G) \to \{0, \ldots, q-1\}$ such that~$f(u) \neq f(v)$ for all~$uv \in E(G)$. A graph~$G$ is \emph{bipartite} if its vertex set can be partitioned into two \emph{partite sets}~$L \dot \cup R$ such that no edge has both of its endpoints in the same partite set. It is well-known that the following three conditions are equivalent for any graph~$G$: (1)~$G$ is bipartite,~(2)~$G$ admits a proper $2$-coloring, and~(3)~there is no cycle of odd length in~$G$. An \emph{odd cycle transversal} (OCT) of a graph~$G$ is a set~$S \subseteq V(G)$ such that~$G - S$ is bipartite. An \emph{independent set} is a vertex set~$S$ such that~$G[S]$ is edgeless. We say that a vertex set~$X$ in a graph~$G$ \emph{separates} two (not necessarily) disjoint vertex sets~$S$ and~$T$ if no connected component of~$G - X$ simultaneously contains a vertex from~$S$ and a vertex from~$T$. For a collection~$\{T_1, \ldots, T_m\}$ of (not necessarily disjoint) vertex sets in a graph~$G$, we say that a vertex set~$X$ is an \emph{$\{T_1, \ldots, T_m\}$-separator} if~$X$ separates all pairs~$(T_i, T_j)$ for~$i \neq j$. Note that~$X$ is allowed to intersect~$\bigcup _{i \in [m]} T_i$.

The next lemma gives a simple sufficient condition for a graph to be bipartite.

\begin{restatable}{lemma}{lemSeparatedColoringsMakeBipartite} \label{lem:separated:colorings:make:bipartite}
	Let $G$ be a graph and let $V_L \cup V_0 \cup V_R = V(G)$ be a partition of its vertices such that $V_0$ is a $\{V_L, V_R\}$-separator. If there exist proper $2$-colorings $f_L: (V_0 \cup V_L) \rightarrow \{0, 1\}$ and $f_R: (V_0 \cup V_R) \rightarrow \{0, 1\}$ of $G[V_0 \cup V_L]$ and $G[V_0 \cup V_R]$ respectively such that $f_L(v_0) = f_R(v_0)$ for every $v_0 \in V_0$, then $G$ is bipartite.
\end{restatable}
\begin{proof}
	To show that $G$ is bipartite, we provide a proper $2$-coloring of the graph. We define this coloring $f \colon V(G) \rightarrow \{0, 1\}$ such that $f(v_0) = f_L(v_0) (= f_R(v_0))$ for every $v_0 \in V_0$, $f(v_L) = f_1(v_L)$ for every $v_L \in V_L$ and $f(v_R) = f_2(v_R)$ for every $v_R \in V_R$. To see that~$f$ is a proper $2$-coloring, we show that no edge~$e \in E(G)$ is monochromatic under~$f$. 
	
	By the assumption that~$V_0$ is a separator, each edge~$e \in E(G)$ is contained in~$G[V_0 \cup V_L]$ or~$G[V_0 \cup V_R]$ (or both). If~$e$ is an edge in the former, its endpoints are colored the same as in~$f_L$ and are therefore bichromatic. The analogous argument for~$f_R$ holds when~$e$ is an edge of the latter.
\end{proof}

The next lemma captures the main idea behind the iterative compression algorithm~\cite{ReedSV04} (cf.~\cite[\S 4.4]{CyganFKLMPPS15}) for solving \textsc{Odd Cycle Transversal}. Given a (potentially suboptimal) odd cycle transversal~$W$ of a graph, it shows that the task of finding an odd cycle transversal disjoint from~$W$ whose removal leaves a bipartite graph with~$W_0, W_1 \subseteq W$ in opposite partite sets of its bipartition is equivalent to separating two vertex sets derived from a baseline bipartition of~$G-W$. Our statement below is implied by Claim 1 in the work of Jansen and de Kroon~\cite{JansenK21}.
\begin{lemma}[{\cite[Claim 1]{JansenK21}}] \label{lem:AR-separation}
	Let~$W$ be an OCT in graph~$G$. For each partition of $W = W_0 \cup W_1$ into two independent sets, for each proper $2$-coloring $c$ of $G - W$, we have the following equivalence for each $X \subseteq V(G) \setminus W$: the graph $G - X$ has a proper $2$-coloring with $W_0$ color $0$ and $W_1$ color $1$ \emph{if and only if} the set $X$ separates $A$ from $R$ in the graph $G - W$, with:
	\begin{align*}
		A &= (N_G(W_0) \cap c^{-1}(0)) \cup (N_G(W_1) \cap c^{-1}(1)), \\ 
		R &= (N_G(W_0) \cap c^{-1}(1)) \cup (N_G(W_1) \cap c^{-1}(0)).
	\end{align*}
\end{lemma}

\subparagraph*{Multiway cuts} Let~$\mathcal{T} = (T_1, \ldots, T_s)$ be a partition of a set~$T \subseteq V(G)$ of \emph{terminal} vertices in an undirected graph~$G$. A \emph{multiway cut} of~$\mathcal{T}$ in~$G$ is a vertex set~$X \subseteq V(G)$ such that for each pair~$t_i, t_j \in T \setminus X$ that belong to different parts of partition~$\mathcal{T}$, the graph~$G-X$ does not contain a path from~$t_i$ to~$t_j$. A \emph{restricted} multiway cut of~$\mathcal{T}$ is a vertex set~$X$ that is a multiway cut for~$\mathcal{T}$ such that~$X \cap T = \emptyset$, i.e., it does not contain any terminals. 

For a positive integer~$s$, a \emph{generalized $s$-partition} of a set~$T$ is a partition~$\mathcal{T}^* = (T_0, T_1, \ldots, T_s, \linebreak[1] T_X)$ of~$T$ into~$s+2$ parts, some of which can be empty. The parts~$T_0$ and~$T_X$ play a special role, which are the \emph{free} and \emph{deleted} part of~$\mathcal{T}^*$, respectively. Let~$T' = T_1 \cup \ldots \cup T_s$. A \emph{multiway cut} of $\mathcal{T}^*$ is a (non-restricted) multiway cut in~$G - T_X$ of the partition $\mathcal{T} = (T_1, \ldots, T_s)$ of~$T'$. Hence the vertices of~$T_X$ are deleted from the graph, while no cut constraints are imposed on the vertices of~$T_0$. 

A \emph{minimum multiway cut} of a generalized $s$-partition $\mathcal{T}^*$ in a graph~$G$ is a minimum-cardinality vertex set that satisfies the requirements of a multiway cut for $\mathcal{T}^*$. We denote the size of a minimum multiway cut of~$\mathcal{T}^*$ in~$G$ by~$\mwc(G, \mathcal{T}^*)$. The following cut covering lemma by Kratsch and Wahlstr\"{o}m will be useful for our algorithm.

\begin{theorem}[{\cite[Theorem 5.14]{KratschW20}}]\label{thm:multiway-cover}
	Let~$G$ be an undirected graph on $n$ vertices with a set~$T \subseteq V(G)$ of terminal vertices, and let~$s \in \mathbb{N}$ be a constant. There is a set~$Z \subseteq V(G)$ with~$|Z| = \Oh(|T|^{s+1})$ such that~$Z$ contains a minimum multiway cut of every generalized $s$-partition~$\mathcal{T}^*$ of~$T$, and we can compute such a set in randomized polynomial time with failure probability~$\mathcal{O}(2^{-n})$.
\end{theorem}

For a generalized $s$-partition~$\mathcal{T} = (T_0, T_1, \ldots, T_s, T_X)$ of a terminal set~$T \subseteq V(G)$ in an undirected graph~$G$, we call a multiway cut~$X$ of~$\mathcal{T}$ \emph{restricted} if it satisfies~$X \cap (\bigcup_{i=1}^s T_i) = \emptyset$. Hence a restricted multiway cut does not delete any vertex that is active as a terminal in the generalized partition. A minimum \emph{restricted} multiway cut of~$\mathcal{T}$ is a restricted multiway cut whose size is minimum among all restricted multiway cuts. We denote the minimum size of a \emph{restricted} multiway cut of~$\mathcal{T}$ in~$G$ by~$\rmwc(G, \mathcal{T})$, which we define as~$+\infty$ if no such cut exists. 

The following lemma shows that the randomization in the polynomial-time algorithm by Kratsch and Wahlstr\"{o}m can be avoided by the use of a single-exponential FPT algorithm, and that the cut covering set can be adapted to work for \emph{restricted} multiway cuts as long as we have a bound on their size.

\begin{restatable}{lemma}{lemCoveringCorollary}[$(\bigstar)$]
	\label{lem:covering:corollary}
	Let~$s \in \mathbb{N}$ be a constant. There is a \emph{deterministic} algorithm that, given an undirected $n$-vertex graph~$G$ and a set~$T \subseteq V(G)$ of terminals, runs in time~$2^{\Oh(|T|)} \cdot n^{\Oh(1)}$ and computes a set~$Z \subseteq V(G)$ with~$|Z| = \Oh(|T|^{2s+2})$ with the following guarantee: for each generalized $s$-partition~$\mathcal{T}$ of~$T$, if there is a restricted multiway cut for~$\mathcal{T}$ of size at most~$|T|$ in~$G$, then the set~$Z$ contains a \emph{minimum} restricted multiway cut of~$\mathcal{T}$.
\end{restatable}